% =====================================================================
%  Математик I В — слайд ба хэвлэх хувилбарт хуваалцах преамбл
%  slides.tex (beamer) ба print.tex (book + beamerarticle) хоёулаа \input хийнэ.
%  \mode<presentation>{...} зөвхөн слайдад, \mode<article>{...} зөвхөн хэвлэхэд.
%  XeLaTeX-ээр хөрвүүлнэ (Overleaf: Menu -> Compiler -> XeLaTeX; latexmkrc үүнийг
%  автоматаар тохируулна).
% =====================================================================

% ---------- Хэл ба фонт ----------
\usepackage{fontspec}
\defaultfontfeatures{Ligatures=TeX}
\usepackage{polyglossia}
\setdefaultlanguage{mongolian}
\setotherlanguage{english}

% Фонтууд TeX Live-д (тиймээс Overleaf-д) байдаг файлуудаас ачаалагдана.
% Слайд: Noto Sans (дэлгэц дээр тод уншигдана); хэвлэх: CMU Serif (Computer Modern-ийн
% кирилл хувилбар, томьёотой нэг загвартай).
\newcommand{\MathSansNoto}{%
  \setsansfont{NotoSans}[Extension=.ttf, UprightFont=*-Regular, BoldFont=*-Bold,
    ItalicFont=*-Italic, BoldItalicFont=*-BoldItalic]%
  \newfontfamily\cyrillicfontsf{NotoSans}[Extension=.ttf, UprightFont=*-Regular,
    BoldFont=*-Bold, ItalicFont=*-Italic, BoldItalicFont=*-BoldItalic]}
\newcommand{\MathSansCMU}{%
  \setsansfont{cmunss.otf}[BoldFont=cmunsx.otf, ItalicFont=cmunsi.otf, BoldItalicFont=cmunso.otf]%
  \newfontfamily\cyrillicfontsf{cmunss.otf}[BoldFont=cmunsx.otf, ItalicFont=cmunsi.otf,
    BoldItalicFont=cmunso.otf]}
\setmainfont{cmunrm.otf}[BoldFont=cmunbx.otf, ItalicFont=cmunti.otf, BoldItalicFont=cmunbi.otf]
\newfontfamily\cyrillicfont{cmunrm.otf}[BoldFont=cmunbx.otf, ItalicFont=cmunti.otf,
  BoldItalicFont=cmunbi.otf]
\setmonofont{cmuntt.otf}[BoldFont=cmuntb.otf, ItalicFont=cmunit.otf, BoldItalicFont=cmuntx.otf]
\newfontfamily\cyrillicfonttt{cmuntt.otf}[BoldFont=cmuntb.otf, ItalicFont=cmunit.otf,
  BoldItalicFont=cmuntx.otf]
\mode<presentation>{\IfFontExistsTF{NotoSans-Regular.ttf}{\MathSansNoto}{\MathSansCMU}}
\mode<article>{\MathSansCMU}

% ---------- Багцууд ----------
\usepackage{amsmath,mathtools,amssymb,amsfonts}
\mode<article>{\usepackage{amsthm}}
\usepackage{mathrsfs}
\usepackage{newunicodechar}
\usepackage{enumitem}
\usepackage{multicol}
\usepackage{environ}
\usepackage{tikz}
\usetikzlibrary{arrows,arrows.meta,calc,patterns,decorations.pathreplacing,positioning}
\usepackage{pgfplots}
\pgfplotsset{compat=1.18}
\usepackage[export]{adjustbox}
\graphicspath{{pictures/}}

% Зураг/хүснэгт хуудсанд (слайдад) багтахгүй бол л жижигрүүлнэ
\mode<presentation>{\newenvironment{fitbox}{\adjustbox{max width=\linewidth,max totalheight=.68\textheight}\bgroup}{\egroup}}
\mode<article>{\newenvironment{fitbox}{\adjustbox{max width=\linewidth,max totalheight=.85\textheight}\bgroup}{\egroup}}
% Слайдад олон баганыг нэг багана болгоно: allowframebreaks зөвхөн нэг баганыг
% хуудас хооронд таслаж чадна. Хэвлэх хувилбарт multicols хэвээр.
\mode<presentation>{\RenewDocumentEnvironment{multicols}{m o o}{}{}}

% Эх файлд шууд бичигдсэн Юникод тэмдэгтүүд
\newunicodechar{→}{\ensuremath{\to}}
\newunicodechar{⇔}{\ensuremath{\Leftrightarrow}}
\newunicodechar{∞}{\ensuremath{\infty}}

% beamer-ийн itemize/enumerate загварыг enumitem эвдэхгүй байлгах
\mode<presentation>{\setlist{label*=\usebeamertemplate{itemize item}}
  \setlist[enumerate]{label=\arabic*.}
  \setlist[enumerate,2]{label=\alph*)}
  \setlist{topsep=2pt,itemsep=2pt,leftmargin=1.5em}}
\mode<article>{\setlist[itemize]{topsep=3pt,itemsep=3pt,leftmargin=1.3em}
  \setlist[enumerate]{topsep=4pt,itemsep=4pt,leftmargin=1.6em}}

% ---------- Макро ----------
\newcommand{\coursename}{Нэг хувьсагчийн функцийн онол 1В (MATH102)}
\newcommand{\semester}{2026 оны намар}
\newcommand{\smallbigcup}{\mathop{\textstyle\bigcup}}
\newcommand{\smallbigcap}{\mathop{\textstyle\bigcap}}
\newcommand{\N}{\mathbb{N}}
\newcommand{\Z}{\mathbb{Z}}
\newcommand{\Q}{\mathbb{Q}}
\newcommand{\R}{\mathbb{R}}
\newcommand{\Sx}{\mathrm{S}}
\newcommand{\itcol}[1]{{\itshape\textcolor{red}{#1}}}
\newcommand{\arccot}{\mathbf{arccot}}
\providecommand{\tg}{\operatorname{tg}}
\providecommand{\ctg}{\operatorname{ctg}}
\providecommand{\arctg}{\operatorname{arctg}}
\providecommand{\arcctg}{\operatorname{arcctg}}

% ---------- Теорем маягийн орчнууд ----------
% beamer-ийн \note (илтгэгчийн тэмдэглэл) командыг "Санамж" орчноор солино
\let\note\relax\let\endnote\relax
\mode<article>{%
  \newtheorem{theorem}{Теорем}[chapter]
  \newtheorem{proposition}{Өгүүлбэр}[chapter]
  \newtheorem{lemma}{Лемма}[section]
  \theoremstyle{definition}
  \newtheorem{definition}{Тодорхойлолт}[chapter]
  \newtheorem{example}{Жишээ}[chapter]
  \theoremstyle{remark}
  \newtheorem*{remark}{Тэмдэглэл}
  \newtheorem*{prof}{Баталгаа}
  \newtheorem*{note}{Санамж}
  \newtheorem*{cor}{Мөрдөлгөө}
  \renewcommand{\proofname}{Баталгаа}}
% Слайдад хайрцаггүй (allowframebreaks хуудас хооронд таслах боломжтой) хувилбар
\mode<presentation>{%
  \newcommand{\MathEnv}[4]{% #1 орчин, #2 нэр, #3 өнгө, #4 дугаартай эсэх (1/0)
    \ifnum#4=1 \newcounter{#1}\fi
    \newenvironment{#1}[1][]{%
      \par\medskip
      \ifnum#4=1 \refstepcounter{#1}\fi
      \begin{trivlist}\item[\hskip\labelsep{\color{#3}\bfseries #2\ifnum#4=1~\thechapter.\arabic{#1}\fi.%
        \if\relax\detokenize{##1}\relax\else\ (##1)\fi}]\ignorespaces}%
      {\end{trivlist}\par\medskip}}
  \newcounter{chapter}
  \MathEnv{theorem}{Теорем}{mathred}{1}
  \MathEnv{proposition}{Өгүүлбэр}{mathred}{1}
  \MathEnv{lemma}{Лемма}{mathred}{1}
  \MathEnv{definition}{Тодорхойлолт}{mathblue}{1}
  \MathEnv{example}{Жишээ}{mathgreen}{1}
  \MathEnv{remark}{Тэмдэглэл}{mathgray}{0}
  \MathEnv{note}{Санамж}{mathgray}{0}
  \MathEnv{cor}{Мөрдөлгөө}{mathred}{0}
  \newenvironment{prof}[1][]{\par\smallskip\begin{trivlist}\item[\hskip\labelsep{\itshape Баталгаа.}]\ignorespaces}%
    {\hfill$\square$\end{trivlist}\par\medskip}
  \providecommand{\qedhere}{}
  \providecommand{\paragraph}{}
  \renewcommand{\paragraph}[1]{\par\medskip\noindent{\bfseries #1}\enskip}
  \let\proof\relax\let\endproof\relax
  \newenvironment{proof}[1][]{\begin{prof}}{\end{prof}}}

% Бодлого ба бодолт
\newcounter{problem}
\newenvironment{problem}[1][]{\refstepcounter{problem}\par\noindent\textbf{Бодлого~\theproblem.}~#1}{\par}
\newif\ifshowsolutions
\showsolutionstrue   % бодолтыг нуух бол \showsolutionsfalse
\NewEnviron{solution}{\ifshowsolutions\par\smallskip\noindent\textit{Бодолт. }\BODY\par\smallskip\fi}


% ---------- Семинар ба гэрийн даалгаврын бодлогууд (Г.Н. Берман) ----------
% \ProblemSet{Семинар 10}  -> шинэ хэсэг, дугаарлалт 1-ээс
% \pgroup{заавар}          -> бүлэг бодлогын заавар
% \begin{plist} \prob{936}{$x=a\cos t$} \end{plist}  -> [Б.936] нь Берманы дугаар
\newcounter{pnum}
\mode<presentation>{\newcommand{\ProblemSet}[1]{\setcounter{pnum}{0}\section{#1}}}
\mode<article>{\newcommand{\ProblemSet}[1]{\setcounter{pnum}{0}\section*{#1}\addcontentsline{toc}{section}{#1}}}
\newcommand{\bn}[1]{\hfill{\color{mathgray}\scriptsize[Б.#1]}}
\newcommand{\prob}[2]{\stepcounter{pnum}\item[\textbf{\thepnum.}] #2\bn{#1}}
\newcommand{\pgroup}[1]{\par\smallskip\noindent\textbf{#1}\par}
\mode<presentation>{\newenvironment{plist}{\begin{enumerate}[leftmargin=2.2em,itemsep=3pt]}{\end{enumerate}}}
\mode<article>{\newenvironment{plist}{\begin{multicols}{2}\begin{enumerate}[leftmargin=2.2em,itemsep=2pt]}{\end{enumerate}\end{multicols}}}

% ---------- Лекцийн гарчиг ----------
% \MathLecture{дугаар}{гарчиг}: слайдад гарчгийн хуудас, хэвлэхэд \chapter
\newcommand{\MathLecture}[2]{%
  \setcounter{chapter}{\numexpr#1-1\relax}%
  \mode<article>{\chapter{#2}}%
  \mode<presentation>{\stepcounter{chapter}%
    \setcounter{definition}{0}\setcounter{theorem}{0}\setcounter{example}{0}%
    \setcounter{proposition}{0}\setcounter{lemma}{0}%
    \title[Лекц #1]{Лекц #1. #2}\subtitle{\coursename}\date{\semester}%
    \begin{frame}[plain]\titlepage\end{frame}%
    \begin{frame}{Агуулга}\tableofcontents\end{frame}}}
